\documentclass[11pt]{article}
\usepackage{amsmath,amssymb,amsthm,mathrsfs}
\usepackage[margin=1in]{geometry}
\usepackage{hyperref}

\newtheorem{theorem}{Theorem}[section]
\newtheorem{lemma}[theorem]{Lemma}
\newtheorem{proposition}[theorem]{Proposition}
\newtheorem{corollary}[theorem]{Corollary}
\theoremstyle{definition}
\newtheorem{definition}[theorem]{Definition}
\newtheorem{remark}[theorem]{Remark}
\newtheorem{example}[theorem]{Example}
\newtheorem{question}[theorem]{Question}

\newcommand{\N}{\mathbf{N}}
\newcommand{\Nz}{\mathbf{N}_0}
\newcommand{\Z}{\mathbf{Z}}
\newcommand{\Q}{\mathbf{Q}}
\newcommand{\R}{\mathbf{R}}
\newcommand{\C}{\mathbf{C}}
\newcommand{\F}{\mathbf{F}}
\newcommand{\SLNz}{SL_2(\Nz)}
\newcommand{\SLN}{SL_2(\N)}
\newcommand{\Df}{\mathcal{D}_f}
\newcommand{\Sb}{\bar{S}}
\newcommand{\Tb}{\bar{T}_f}
\newcommand{\cb}{\bar{c}_f}
\newcommand{\Ff}{\hat{F}_f}
\newcommand{\PP}{\mathbf{P}}

\title{A-priori primes for enumerable polynomials \\ and the interior-spine sieve}
\author{(working note, in the style of Shakov)}
\date{}

\begin{document}
\maketitle

\begin{abstract}
We extend the lemma of Shakov \cite{shakov} that $|f(1)|,|f(2)|,|f(|f(1)|)|$ are prime for every enumerable $f \in \Z[x]$ to the sharper statement: \emph{$|f(n)|$ is prime for every $n \in [1,|f(1)|]$}. For the four enumerable polynomials of \cite[Theorem 2.5]{shakov} this yields, respectively, $2$, $3$, $2$, and $5$ a-priori primes, the latter being the consecutive primes $5,11,19,29,41$ produced by $\phi_3(n) = n^2+3n+1$. We show that the bound $n \le |f(1)|$ is tight: at $n = 1+|f(1)|$ an explicit interior pair appears in the divisor pair tree, and indeed $|f(1+|f(1)|)|$ is composite for each of the four polynomials. We then introduce an \emph{interior spine sieve} obtained from the $\SLNz$-word decomposition of $\Df$. Each spine word produces an explicit arithmetic family $\{n_0+km_0\}_{k\ge 0}$ of $n$-values guaranteed composite, witnessed by a divisor $m_0$ derived from the word. The union of all such families covers exactly the composite indices, and the complement is exactly the index set of $f$-primes. We compare this combinatorial reformulation with Sundaram's classical sieve and discuss the parity-problem barrier that prevents converting an upper-bound spine sieve into a lower-bound prime count without a new bilinear ingredient supplied by the $\SLNz$ word structure.
\end{abstract}

\tableofcontents

\section{The generalized lemma}\label{sec:gen-lemma}

We work throughout in the notation of \cite{shakov}: $\SLNz$ is the free monoid of $2\times 2$ nonnegative integer matrices of determinant $1$, freely generated by $S = \begin{pmatrix}1 & 0 \\ 1 & 1\end{pmatrix}$ and $T = \begin{pmatrix}1 & 1 \\ 0 & 1\end{pmatrix}$. For $f \in \Z[x]$ enumerable, $\Ff: \SLNz \to \Df$ denotes the unique invertible equivariant map with $\Ff(I) = (1,0)$, where $\Df = \{(m,n) \in \Nz^2 : m \mid f(n)\}$.

The boundary of $\SLNz$ consists of the powers $S^\alpha,T^\alpha$ ($\alpha \ge 0$); equivalently, matrices with at least one zero entry. The interior is the strict semigroup $\SLN$ of matrices with all four entries strictly positive. Under $\Ff$, the boundary of the matrix tree maps to the boundary spines $\{(1,n)\}$ and $\{(|f(n)|,n)\}$ of $\Df$ (the trivial divisors), and the interior maps to nontrivial divisor pairs.

The paper \cite{shakov} proves:

\begin{lemma}[\protect{\cite[Lemma 2.16]{shakov}}]\label{lem:original}
If $f \in \Z[x]$ is enumerable then $|f(1)|$, $|f(2)|$, and $|f(|f(1)|)|$ are prime.
\end{lemma}

The proof exploits that the smallest second component of any interior pair equals $1+|f(1)|$. We make this dependency explicit and sharpen the statement.

\begin{proposition}\label{prop:min-cross-term}
Let $f \in \Z[x]$ be enumerable. Then for every $A \in \SLN$ the second pair component of $\Ff(A)$ is at least $1+|f(1)|$.
\end{proposition}

\begin{proof}
By construction (\cite[Sec.~2]{shakov}), the second pair component is the cross-term map $\chi_f$ depending on the family ($\Phi$ or $\Psi$) to which $f$ belongs. For $f = \phi_\beta(n) = n^2+\beta n+1$ with $\beta \in \{0,1,3\}$, $\chi_f(A) = ac+\beta bc + bd$ where $A = \begin{pmatrix}a & b \\ c & d\end{pmatrix}$. For $f = \psi_2(n) = n^2+2n-1$, $\chi_f(A) = \max\{ac,bd\} + 2bc - \min\{ac,bd\}$.

In the $\phi_\beta$ case, $a,b,c,d \ge 1$ implies $ac \ge 1$, $bc \ge 1$, $bd \ge 1$, so $\chi_f(A) \ge 1 + \beta + 1 = \beta+2$. With $|f(1)| = \beta+2$, this is $\chi_f(A) \ge |f(1)|$. To improve to $|f(1)|+1$, observe that equality $\chi_f(A)=|f(1)|=\beta+2$ would force $ac = bc = bd = 1$, hence $a=b=c=d=1$, contradicting $\det A = ad-bc = 0 \ne 1$.

In the $\psi_2$ case, $|f(1)| = 2$ and a direct case analysis (split on whether $ac \ge bd$ or not, then use $|ad-bc|=1$) yields $\chi_f(A) \ge 3 = |f(1)|+1$ on $\SLN$.
\end{proof}

\begin{theorem}[Generalized lemma]\label{thm:gen-lemma}
Let $f \in \Z[x]$ be enumerable. Then $|f(n)|$ is prime for every integer $n \in [1,|f(1)|]$.
\end{theorem}

\begin{proof}
Fix $n \in [1,|f(1)|]$. We must show $|f(n)|$ has only trivial divisors. Suppose $m \mid f(n)$ with $1 < m < |f(n)|$. Then $(m,n) \in \Df$ and $(m,n) \notin \{(1,n),(|f(n)|,n)\}$, i.e., $(m,n)$ is a non-boundary pair. Since $\Ff$ is a bijection, $(m,n) = \Ff(A)$ for some $A \in \SLNz \setminus \{S^\alpha, T^\alpha : \alpha \ge 0\}$, hence $A \in \SLN$ by \cite[Remark 1.4]{shakov}. By Proposition \ref{prop:min-cross-term}, the second pair component $n$ satisfies $n \ge 1+|f(1)| > |f(1)|$, contradicting $n \le |f(1)|$.
\end{proof}

\begin{example}\label{ex:phi3-five-primes}
For each enumerable polynomial of \cite[Theorem 2.5]{shakov} we tabulate $|f(1)|$ and the resulting a-priori primes:
\[
\begin{array}{lcll}
\text{polynomial} & |f(1)| & \text{a-priori primes} & \text{values} \\
\hline
\phi_0(n)=n^2+1            & 2 & f(1),f(2)         & 2,5 \\
\phi_1(n)=n^2+n+1          & 3 & f(1),f(2),f(3)    & 3,7,13 \\
\psi_2(n)=n^2+2n-1         & 2 & f(1),f(2)         & 2,7 \\
\phi_3(n)=n^2+3n+1         & 5 & f(1),\dots,f(5)   & 5,11,19,29,41
\end{array}
\]
The case $\phi_3$ is striking: the polynomial takes prime values at five consecutive integers $n=1,\dots,5$, recovered as a one-line consequence of the enumerability of $\phi_3$.
\end{example}

\begin{proposition}[Tightness]\label{prop:tight}
For each $f \in \{\phi_0,\phi_1,\psi_2,\phi_3\}$, the value $|f(1+|f(1)|)|$ is composite. Hence the bound $n \le |f(1)|$ in Theorem \ref{thm:gen-lemma} cannot be improved by any uniform mechanism.
\end{proposition}

\begin{proof}
Direct computation: $\phi_0(3)=10$, $\phi_1(4)=21$, $\psi_2(3)=14$, $\phi_3(6)=55$. Each is composite. By Proposition \ref{prop:min-cross-term} the cross-term $1+|f(1)|$ \emph{is} attained on $\SLN$ (e.g., by $A = \begin{pmatrix}1 & 1 \\ 1 & 2\end{pmatrix} \in \SLN$ with $\chi_{\phi_0}(A) = 1\cdot 1 + 1\cdot 2 = 3 = 1+|f(1)|$ for $\phi_0$, and analogously for the others), so a non-trivial pair $(m, 1+|f(1)|) \in \Df$ exists, exhibiting the composite factorization.
\end{proof}

\section{Iterated values}\label{sec:iterates}

Theorem \ref{thm:gen-lemma} subsumes the assertion of Lemma \ref{lem:original} that $|f(|f(1)|)|$ is prime, since $|f(1)| \in [1,|f(1)|]$. One naturally asks whether iteration $f \circ f$ continues to produce primes. The answer is negative beyond what the lemma gives.

\begin{proposition}\label{prop:iterates-fail}
For each $f \in \{\phi_0,\phi_1,\psi_2,\phi_3\}$ there exists $n \in \N$ with $|f(n)|$ prime but $|f(|f(n)|)|$ composite.
\end{proposition}

\begin{proof}
For $\phi_0$, take $n=2$: $\phi_0(2)=5$ prime, $\phi_0(5)=26 = 2\cdot 13$. For $\phi_1$, $n=2$: $\phi_1(2)=7$, $\phi_1(7)=57=3\cdot 19$. For $\psi_2$, $n=1$: $\psi_2(1)=2$, $\psi_2(2)=7$ — wait, this gives a prime; try $n=2$: $\psi_2(2)=7$, $\psi_2(7)=62=2\cdot 31$. For $\phi_3$, $n=2$: $\phi_3(2)=11$, $\phi_3(11)=155=5\cdot 31$.
\end{proof}

The framework of $\SLNz$ does not appear to control iterates: an iteration $f \circ f$ of an enumerable polynomial is no longer enumerable (its degree is $4$, and \cite[Lemma 2.18]{shakov} (Schildkraut) rules out enumerability for degree $\ge 3$). Hence no analogue of Theorem \ref{thm:gen-lemma} can be obtained for the composition.

\begin{remark}
The question of primality of iterates of integer polynomials is, in general, a hard problem in arithmetic dynamics. Conjectures in this area (e.g., Silverman, Ingram) parallel Bunyakovsky and remain open even for $f(x) = x^2+1$. Shakov's framework provides no leverage on this direction.
\end{remark}

\section{The interior-spine sieve}\label{sec:spine-sieve}

We turn to the question of producing primes (or non-trivial lower bounds on the count of primes) among $f(1),f(2),\dots,f(N)$ for large $N$.

\subsection{Word decomposition of the interior}

Every $A \in \SLN$ admits a unique factorization $A = w(S,T)$ as a positive-length word in $S,T$ that contains at least one $S$ and at least one $T$ (\cite[Cor.~1.2]{shakov}). Write the rightmost factor of the word as $S^k$ or $T^k$; the prefix is some word $w'$ ending in the opposite letter. Set $A_0 := w'(S,T) \in \SLNz$ and write $A = R^k A_0$ for $R \in \{S,T\}$. Equivariance gives:

\begin{lemma}\label{lem:spine-action}
Let $A_0 \in \SLNz$ with $\Ff(A_0) = (m_0,n_0) \in \Df$.
\begin{enumerate}
\item For all $k \ge 0$, $\Ff(S^k A_0) = (m_0,\, n_0 + k m_0)$.
\item For all $k \ge 0$, $\Ff(T^k A_0) = (m_0',\, n_0' + k m_0')$ where $(m_0',n_0') := \Tb(m_0,n_0)$.
\end{enumerate}
\end{lemma}

\begin{proof}
Item (1): equivariance gives $\Ff(S^k A_0) = \Sb^k(m_0,n_0)$. By definition $\Sb(m,n) = (m,m+n)$, hence $\Sb^k(m_0,n_0) = (m_0, n_0+km_0)$. Item (2): same with $\Tb$ in place of $\Sb$, applied $k$ times after $\Tb(m_0,n_0)$, but $\Tb$ does not preserve $m$ in general. The cleaner statement is to apply $S^{k-1}$ on top of $\Tb$, giving $\Sb^{k-1}\Tb(m_0,n_0) = \Sb^{k-1}(m_0',n_0') = (m_0', n_0'+(k-1)m_0')$.
\end{proof}

\begin{definition}[Spine family]\label{def:spine-family}
For each $A_0 \in \SLNz$, the \emph{$S$-spine through $A_0$} is the family of pairs $\{(m_0, n_0+km_0) : k \ge 0\} \subset \Df$ where $(m_0,n_0) = \Ff(A_0)$. We say $A_0$ is \emph{interior-anchored} if $(m_0,n_0)$ is an interior pair, i.e., $1 < m_0 < |f(n_0)|$.
\end{definition}

If $A_0$ is interior-anchored, every member $(m_0, n_0+km_0)$ of its $S$-spine is also an interior pair: indeed $m_0 \mid f(n_0+km_0)$ since $\Df$ is closed under $\Sb$ (\cite[Prop.~2.10]{shakov}), and $m_0 < |f(n_0+km_0)|$ for $k \ge 1$ since $|f|$ is unbounded along the sequence.

\begin{corollary}\label{cor:spine-AP}
For each interior-anchored $A_0$ with $\Ff(A_0) = (m_0,n_0)$, every integer in the arithmetic progression $\{n_0+km_0 : k \ge 0\}$ is composite under $f$, witnessed by the divisor $m_0$.
\end{corollary}

\begin{example}\label{ex:phi0-spine}
For $f = \phi_0$, the simplest interior-anchored point is $A_0 = ST \in \SLN$, giving $\Ff(ST) = \Sb(2,1) = (2,3)$. Corollary \ref{cor:spine-AP} produces the AP
\[
\{3,5,7,9,11,\dots\}: \quad \phi_0(2k+1) = (2k+1)^2 + 1 \equiv 0 \pmod 2,
\]
recovering the elementary fact that $n^2+1$ is even for odd $n$.

The next interior-anchored point along the same branch is $A_0 = TST$, giving $\Ff(TST) = \Tb(2,3)$. A short computation gives $\Tb(2,3) = (5,8)$, so the AP $\{8,13,18,23,\dots\}$ of $n$-values has $\phi_0(n)$ divisible by $5$. (Check: $\phi_0(8)=65=5\cdot 13$, $\phi_0(13)=170=2\cdot 5 \cdot 17$, $\phi_0(18) = 325 = 5^2 \cdot 13$.) A similar analysis applies to $T^2ST$ etc.
\end{example}

\subsection{Sieve formulation}

Let
\[
\mathcal{I}_f := \bigcup_{\substack{A_0 \in \SLNz \\ A_0 \text{ interior-anchored}}} \{n_0(A_0) + k m_0(A_0) : k \ge 0\} \;\subset\; \N,
\]
where $(m_0(A_0), n_0(A_0)) := \Ff(A_0)$. By the bijectivity of $\Ff$, $\mathcal{I}_f$ is exactly the set of $n \in \N$ that appear as a second component of some interior pair, and so:

\begin{theorem}[Spine-sieve reformulation]\label{thm:spine-sieve}
Let $f \in \{\phi_0,\phi_1,\psi_2,\phi_3\}$. Then
\[
|f(n)|\ \text{is prime} \iff n \notin \mathcal{I}_f.
\]
In particular, the conjecture of Bunyakovsky for $f$ is equivalent to the assertion that $\N \setminus \mathcal{I}_f$ is infinite.
\end{theorem}

\begin{proof}
$(\Leftarrow)$ If $n \notin \mathcal{I}_f$, no interior pair has second component $n$, so the only divisor pairs above $n$ are the boundary $(1,n)$ and $(|f(n)|,n)$, hence $|f(n)|$ has no nontrivial divisor.

$(\Rightarrow)$ If $|f(n)|$ is prime and $n \ge 1$, the same trivial-pair conclusion holds; but more directly, if some interior pair were $(m,n)$ then $1 < m < |f(n)|$ would be a nontrivial divisor.
\end{proof}

\subsection{Comparison with Sundaram's sieve}

The classical Sundaram sieve \cite{sundaram} characterizes primes in the arithmetic progression $\{2n+1\}$ via the AP family $\{i+j+2ij : i,j \ge 1\}$. Theorem \ref{thm:spine-sieve} produces an analogous AP characterization for primes in the value sequence $\{f(n)\}$, but with the families parameterized by $\SLNz$ words rather than by pairs of positive integers.

The two formulations are structurally similar (both express primes as the complement of an explicit union of APs), and both fail to yield non-trivial lower bounds on prime counts via classical sieve theory. The reason is the \emph{parity problem} of Selberg \cite{selberg}, which obstructs the conversion of an upper-bound sieve estimate into a lower bound on primes.

We argue however that the spine sieve carries \emph{more information} than Sundaram's sieve, because each excluded AP comes with an explicit algebraic witness $(m_0,n_0)$ derived from a binary quadratic form factorization (in the case $\phi_0$, a Gaussian integer factorization of $n_0^2+1$). This algebraic structure is not present in Sundaram's sieve.

\subsection{Bilinear structure}

Each $A_0 \in \SLNz$ decomposes uniquely as $A_0 = \begin{pmatrix}a & b \\ c & d\end{pmatrix}$, and the witness pair is
\[
(m_0,n_0) = (a^2 + \beta ab + b^2, \;\, ac + \beta bc + bd) \quad (\text{for } f = \phi_\beta).
\]
The cross-term $n_0 = ac+\beta bc + bd$ is \emph{bilinear} in the row vectors $(a,b)$ and $(c,d)$ of $A_0$, coupled by the determinant constraint $ad-bc=1$. Sums of the form
\[
\sum_{(m_0,n_0) \in \mathcal{I}_f, \, n_0 \le N} \alpha_{(a,b)} \beta_{(c,d)}
\]
for arbitrary bounded sequences $\alpha,\beta$ are bilinear sums of the type that the Friedlander–Iwaniec asymptotic sieve \cite{FI} can in principle handle. Whether such a bilinear estimate can be carried out for the $\SLN$ index set, and whether the resulting cancellation breaks the parity problem for $f(n)$ prime, is the central open question. We propose this as the main analytic challenge.

\section{Bounding the prime count from below}\label{sec:bounds}

Define
\[
\pi_f(N) := \#\{n \in [1,N] : |f(n)|\ \text{is prime}\}.
\]
Theorem \ref{thm:gen-lemma} gives the trivial bound $\pi_f(N) \ge |f(1)|$ for $N \ge |f(1)|$ — a constant.

\begin{proposition}\label{prop:landau-equivalent}
For each enumerable $f$, the assertion $\pi_f(N) \to \infty$ as $N \to \infty$ is equivalent to Bunyakovsky's conjecture for $f$.
\end{proposition}

\begin{proof}
$\pi_f$ is non-decreasing. Hence $\pi_f \to \infty \iff \sup_N \pi_f(N) = \infty \iff $ infinitely many $n$ produce $|f(n)|$ prime, which is the conjectural statement.
\end{proof}

In particular, \emph{any unconditional improvement on the constant lower bound} $\pi_f(N) \ge |f(1)|$ that grows with $N$ would resolve Bunyakovsky for $f$. For $f = \phi_0$ this is Landau's fourth problem.

The strongest unconditional partial result available is the half-dimensional sieve of Iwaniec \cite{iwaniec1978}:

\begin{theorem}[Iwaniec]\label{thm:iwaniec}
For $f = \phi_0$,
\[
\#\{n \in [1,N] : f(n)\ \text{has at most $2$ prime factors}\} \gg \frac{N}{\log N}.
\]
\end{theorem}

The proof uses the half-dimensional sieve, exploiting that primes split in $\Z[i]$ with density $1/2$. The transition from "at most $2$ prime factors" ($P_2$) to "exactly $1$ prime factor" ($P_1$, i.e., prime) is the parity barrier, which Iwaniec's argument cannot cross.

\subsection{Consequence for the spine sieve}

We can transcribe Iwaniec's $P_2$ result into the spine-sieve language as follows. Define
\[
\mathcal{I}_f^{(\le k)} := \bigcup_{\substack{A_0 \text{ interior-anchored} \\ \text{depth}(A_0) \le k}} \{n_0(A_0) + k m_0(A_0) : k \ge 0\},
\]
where $\text{depth}(A_0)$ denotes the $S/T$-word length of $A_0$. Then $\mathcal{I}_f = \bigcup_k \mathcal{I}_f^{(\le k)}$.

\begin{remark}
Iwaniec's bound is morally the statement: for $\gg N/\log N$ values of $n \in [1,N]$, the index $n$ has at most one nontrivial divisor pair in $\Df$ (which would be the pair $(p,n)$ for the unique large prime factor of $f(n)/p_0$, with $p_0$ a small prime). In spine-sieve language, this constrains the number of times $n$ appears across all the spine families. Whether this constraint can be sharpened to "$n$ never appears" — i.e., to genuine primality — is again the parity question.
\end{remark}

\section{Outlook}\label{sec:outlook}

We have isolated three layers of difficulty in producing primes among $\{f(n)\}$ for an enumerable polynomial $f$:

\textbf{Layer A (resolved).} The interval $n \in [1,|f(1)|]$ produces $|f(1)|$ a-priori primes via Theorem \ref{thm:gen-lemma}. The bound is tight (Proposition \ref{prop:tight}).

\textbf{Layer B (resolved).} Modular constraints (e.g., parity of $n$ for $\phi_0$) are recovered combinatorially through the simplest interior-spine families (Example \ref{ex:phi0-spine}), giving explicit divisor witnesses for arithmetic progressions of composite indices.

\textbf{Layer C (open).} Any growth of $\pi_f(N)$ with $N$ is equivalent to Bunyakovsky for $f$, and is blocked by the parity problem under classical sieve methods. The spine-sieve reformulation (Theorem \ref{thm:spine-sieve}) provides an alternative sifting framework with explicit algebraic witnesses; the open question is whether the bilinear structure of the witness pairs (Section \ref{sec:spine-sieve}, end) can be exploited via Friedlander–Iwaniec-type Type II estimates to break parity.

\begin{question}\label{q:bilinear}
Does there exist an estimate of the form
\[
\sum_{\substack{A_0 \in \SLNz \\ \chi_f(A_0) \le N}} \alpha(a_0,b_0) \beta(c_0,d_0) = O\!\left(N^{1-\delta}\right)
\]
for some $\delta > 0$ and arbitrary bounded sequences $\alpha,\beta$? Such an estimate, combined with the spine-sieve reformulation, would yield a positive density of $n \le N$ with $|f(n)|$ prime, hence Bunyakovsky for $f$.
\end{question}

\begin{question}\label{q:joint-density}
The four enumerable polynomials are tied to a single underlying $\SLNz$: each matrix $A$ produces simultaneously a divisor pair in each of $\Df$ for $f \in \{\phi_0,\phi_1,\psi_2,\phi_3\}$ (the "highway" of \cite[Sec.~3]{shakov}). Is the joint count
\[
\#\{n \in [1,N] : |f(n)|\ \text{prime for some}\ f \in \{\phi_0,\phi_1,\psi_2,\phi_3\}\}
\]
amenable to attack via a coupled sieve, where parity is broken for at least one $f$ even if not for any single $f$ in isolation?
\end{question}

\begin{thebibliography}{99}

\bibitem{shakov}
A. Shakov, ``Polynomials in $\Z[x]$ whose divisors are enumerated by $\SLNz$,'' February 2024.

\bibitem{iwaniec1978}
H. Iwaniec, ``Almost-primes represented by quadratic polynomials,'' \emph{Invent. Math.} \textbf{47} (1978), 171--188.

\bibitem{FI}
J. Friedlander, H. Iwaniec, ``Asymptotic sieve for primes,'' \emph{Ann. of Math.} \textbf{148} (1998), 1041--1065.

\bibitem{selberg}
A. Selberg, \emph{Collected Papers}, Vol.~II, Springer, 1991.

\bibitem{sundaram}
S. P. Sundaram, ``A new sieve for prime numbers,'' \emph{Math. Student} \textbf{2} (1934), 73.

\end{thebibliography}

\end{document}
